\documentclass[10pt,a4paper]{amsart}
\usepackage[margin=19mm]{geometry}
\usepackage{amsmath,amssymb,mathtools}
\usepackage[scr=rsfs,cal=euler]{mathalfa}
\usepackage{xfrac}
\usepackage[unicode,hidelinks,bookmarks=false]{hyperref}
\newcommand{\ie}{i.e.\ }
\newcommand{\cf}{cf.\ }
\newcommand{\resp}{respectively\ }
\newcommand{\Subsection}{\subsection}
\providecommand{\nobreakdash}{\nobreak\hbox{-}}
\setlength{\emergencystretch}{2em}
\multlinegap0pt
\usepackage{bm}
\usepackage{bbm}
\usepackage{dsfont}
\usepackage{xspace}
\usepackage{cancel}
\usepackage{mathtools}
\usepackage{amsthm}
\makeatletter
\@ifundefined{theorem}{\newtheorem{theorem}{Theorem}}{}
\@ifundefined{lem}{\newtheorem{lem}[theorem]{Lemma}}{}
\makeatother
\title[Holes in Calabi-Yau Effective Cones]{Holes in Calabi-Yau Effective Cones}
\author{Naomi Gendler, Elijah Sheridan, Michael Stillman, David H. Wu}
\date{Source: arXiv:2603.11173v1 (CC BY 4.0)}
\begin{document}
\maketitle

\noindent\emph{Source and licence.} Holes in Calabi-Yau Effective Cones. Version arXiv:2603.11173v1 (CC BY 4.0), distributed under the Creative Commons Attribution 4.0 International licence, as recorded in the arXiv licence metadata for this exact version and shown on the abstract page https://arxiv.org/abs/2603.11173v1. Full rights notice: licenses/arxiv-2603.11173-CC-BY-4.0.txt.\par\smallskip

\noindent\textit{Editorial scope.} Counted unit: the Lem whose source label is \texttt{lem:lift\_to\_global} in the article (source0.tex lines 745--748, source label \texttt{lem:lift\_to\_global}), delivered together with the article's own complete original proof of it (source0.tex lines 749--766, one \texttt{proof} environment of 18 lines). No mathematics is written, repaired or re-derived: statement and proof are the article's own text line for line. Required context retained uncounted: none. Every definition, display and label of those lines is retained unchanged. Omitted from this package and not redistributed: 1--744, 767--1860. Nothing inside a retained excerpt is omitted or altered: every retained body is delivered exactly as the article's lines are written, so no omission receipt of any kind accompanies this package. Citations: the retained excerpts carry no citation command at all, so no bibliography entry of the article is transcribed and no reference is redistributed. Reference adaptation: none was needed and none was made; every reference inside the delivered unit resolves inside it, so no unresolved reference or citation is produced. Printed slips kept literal: no printed word, symbol, hypothesis or number of the counted statement or of its proof was altered. Layout and class: the article's own class is \texttt{article} with packages including setspace, jheppub, inputenc, epsfig, bbm, bbding, amsfonts, amssymb, mathtools, dsfont, bm, graphicx, longtable, pdflscape; the wrapper is the lane's pinned \texttt{amsart} editorial wrapper at 10pt on A4, which restates the theorem-like environments the retained text needs and adds the article's own macro definitions that the retained text needs (none), with extra wrapper packages (bm, bbm, dsfont, xspace, cancel, mathtools). The delivered numbering is the unit's own. Assets: no retained span contains a figure environment, a graphics-inclusion command, a PDF-page inclusion, a table or a TikZ picture; no image file is copied into this package, the package declares no asset dependency and no image file is redistributed with it. Licence: version arXiv:2603.11173v1 of this article is distributed under the Creative Commons Attribution 4.0 International licence, as recorded in the arXiv licence metadata for this exact version and shown on the abstract page https://arxiv.org/abs/2603.11173v1. A full rights notice is delivered at licenses/arxiv-2603.11173-CC-BY-4.0.txt. Characters: every retained excerpt is pure ASCII, so the delivered engine, XeLaTeX with its default Latin Modern text font, reports no missing character.
\begin{lem}
    \label{lem:lift_to_global}
    Let $X$ be a normal projective variety with a Weil divisor $D$ satisfying $[D] \in \mathcal{E}_X$. If $U \subset X$ is an open subset such that the prime codimension-one components of $X \setminus U$ are contained in the stable base locus of $[D]$, then $H^0(X, \mathcal{O}_X(D)) \cong \mathcal{O}_X(D)(U)$.
\end{lem}

\begin{proof}
    Let $E_1, \dots, E_n$ denote the prime codimension-one components of $X \setminus U$. Certainly $H^0(X, \mathcal{O}_X(D)) \subset \mathcal{O}_X(D)(U)$, as any rational function $f$ obeying $\mathrm{div}(f) + D \geq 0$ on all prime divisors will do so for all prime divisors on $U$. To show the other inclusion, we will assume by way of contradiction that $\psi \in \mathcal{O}_X(D)(U) \setminus H^0(X, \mathcal{O}_X(D))$ and use $\psi$ to construct a global section of some multiple of $[D]$ which does not contain some $E_j$ in its base locus. By construction, the divisor $D_\psi := \mathrm{div}(\psi) + D$ has a pole, but it can only be on one of the $E_i$.

    Choose $m > 0$ such that $m[D]$ is effective, and let $D_\varphi := \mathrm{div}(\varphi) + mD \geq 0$ (i.e., $\varphi \in H^0(X, \mathcal{O}_X(mD))$). It will suffice to choose $k,\ell \in \mathbb{Z}_+$ such that the divisor 
    \begin{equation}
        A = k(mD_\psi) + \ell(D_\varphi) = \mathrm{div}(\psi^{mk}\varphi^{\ell}) + m(k+\ell)D
    \end{equation}
    with class $m(k+\ell)[D]$ is effective and does not contain some $E_j$. Note that because $A$ is the sum of a multiple of $D_\varphi$ and an effective divisor, it too can only have poles on the $E_i$. 
    
    To this end, define a function $g : [0,1] \to \mathbb{R}$ by\footnote{One can naturally think of this as the smallest degree of vanishing attained by the $\mathbb{R}$-divisor 
    \begin{equation}
        (1 - t)(m\mathrm{div}(\psi) + mD) + t(\mathrm{div}(\varphi) + mD) = (1 - t)m\mathrm{div}(\psi) + t\mathrm{div}(\varphi) + mD
    \end{equation} across all of the $E_i$, as a function of $t$.}
    \begin{equation}
        g(t) = \mathrm{min}_i\Big((1-t)\nu_{E_i}(mD_\psi) + t \, \nu_{E_i}(\mathrm{div}(D_\varphi)\Big).
    \end{equation}
    We note that $g(0) < 0$, because $\psi^m$ has a pole on some $E_i$. Additionally, $g(1) > 0$, as each $E_i$ belongs to the base locus of $mD$, so the degree of vanishing of $D_\varphi$ on each $E_i$ is at least one. Thus, there is a unique rational $s \in (0,1)$ such that $g(s) = 0$. Pick an integer $c$ that clears denominators such that $k = c(1-s), \ell = cs$ are integral. For these choices of $k, \ell$, the divisor $A$ defined above is effective and doesn't contain some $E_j$. Indeed, as noted, for effectiveness we only need to check for poles on the $E_i$, and $\nu_{E_i}(A) \geq \mathrm{min}_i(\nu_{E_i}(A)) = c g(s) = 0$, meaning $A \geq 0$ but also that $\nu_{E_j}(A) = 0$ for some $j$.
\end{proof}

\end{document}
